\scanpage{52}
Soit de \unsure{m\textsuperscript{e}|\,$=$|même} \deleted{$Z$} $Z$ un espace topologique avec \deleted{\unsure{muni de}} groupe \inserted{discret $G$} qui opère, d'où $\deleted{\unsure{B|\operatorname{Top}}}\operatorname{Top}(Z,G)$, si $H$ est un ss-groupe de $G$, on a un foncteur ``restriction des op. à $H$''

\[\deleted{B}\operatorname{Top}(Z,\unsure{G|C})\longrightarrow\deleted{\unsure{B}}\operatorname{Top}(Z,H)\]

qui s'identifie à un foncteur de localisation, relativement à l'objet \inserted{$Z'$} suivant : posant $X=G/H$, on prend $Z'=Z\times X$ avec opérations diagonales de $G$ et on considère la section \inserted{$e$} de $Z'$ sur $Z$ (\unsure{dans|ds.|c.-à-d.} $\operatorname{Top}(\unsure{Z|z|\mathbb Z})$) définie par l'élémt marqué $x$ de $X$. Ceci dit, on \deleted{prend} le foncteur

\[\begin{aligned}\operatorname{Top}(Z,G)\deleted{\unsure{/}}_{/Z'}&\longrightarrow\operatorname{Top}(Z,H)\\E&\longmapsto E\times_{Z'}(X,e)\end{aligned}\]

est une équivalence de catégories .... (En fait, $Z'$ est un ``revêtement topologique'' de $Z$ ...... )
% archive sheet number: 51
